\documentclass{article}

\usepackage[french]{babel}
\usepackage[utf8]{inputenc}

%%%%%%%%%%%%%%%% Lengths %%%%%%%%%%%%%%%%
\setlength{\textwidth}{15.5cm}
\setlength{\evensidemargin}{0.5cm}
\setlength{\oddsidemargin}{0.5cm}

%%%%%%%%%%%%%%%% Variables %%%%%%%%%%%%%%%%
\def\projet{3}
\def\titre{Interpolation and integration methods / Cubic splines and surface interpolation}
\def\groupe{1}
\def\equipe{2}
\def\responsible{Rafiq Amrar}
\def\secretary{Salah-Eddine Mabkhout}
\def\others{Anas Benbaki, Manal El Abboudi}

\begin{document}

%%%%%%%%%%%%%%%% Header %%%%%%%%%%%%%%%%
\noindent\begin{minipage}{0.98\textwidth}
  \vskip 0mm
  \noindent
  { \begin{tabular}{p{7.5cm}}
      {\bfseries \sffamily
        Projet \projet} \\ 
      {\itshape \titre}
    \end{tabular}}
  \hfill 
  \fbox{\begin{tabular}{l}
      {~\hfill \bfseries \sffamily Groupe \groupe\ - Equipe \equipe
        \hfill~} \\[2mm] 
      Responsable : \responsible \\
      Secrétaire : \secretary \\
      Codeurs : \others
    \end{tabular}}
  \vskip 4mm ~

  ~~~\parbox{0.95\textwidth}{\small \textit{Résumé~:} \sffamily Ce projet porte sur la mise en œuvre de 
  méthodes numériques d’interpolation et d’intégration. Il comporte deux étapes préliminaires : 
  l’interpolation de fonctions à l’aide de splines cubiques, puis l’implémentation de méthodes d’intégration 
  pour évaluer des fonctions données. Ces outils seront ensuite appliqués à l’étude d’un profil d’aile, 
  afin de modéliser l’écoulement de l’air et construire une carte de pression autour de l’extrados et de 
  l’intrados. }
  \vskip 1mm ~
\end{minipage}

%%%%%%%%%%%%%%%% Main part %%%%%%%%%%%%%%%%
\section{Airfoil refinement}
\label{sec:airfoil}

\section{Computing the length of plane curves}
\label{sec:curves}

\section{Modelling the airflow / Pressure map }
\label{sec:pressure}


\end{document}